\chapter{Conclusion and Future Work}

\section{Conclusion}

This project delivers an auditable comparison of Flow Matching, Flow Matching
with logit-normal time sampling, and Mean Flow for unconditional CIFAR-10 image
generation. The implementations share a common backbone, data pipeline,
evaluation path, and evidence format, making configuration differences visible
rather than hiding them inside separate codebases.

Under the submitted training artefacts, FM-LN is the strongest multi-step
configuration at NFE=20, reaching FID~72.66 and IS~$6.18\!\pm\!0.52$ on
1{,}000 generated images. Standard FM is close, with FID~76.10 and
IS~$6.12\!\pm\!0.72$. Mean Flow improves the one-step result relative to both
FM variants, demonstrating its direct displacement mechanism, but its rising
training loss and degraded multi-step behaviour show that the learned average
velocity field is not yet competitive.

These results support a narrow conclusion: logit-normal time sampling is the
best observed configuration among the submitted artefacts, while Mean Flow's
potential depends on substantially more stable interval-conditioned training.
Because training budgets, hardware contexts, and evaluation sample counts are
not fully matched, the evidence does not establish universal algorithm
superiority.

\section{Future Direction}

The project identifies two specific future research contributions, matching the
approved presentation scope.

\subsection{Adaptive MeanFlow for Dynamic NFE}

Instead of assigning every sample the same inference budget, an adaptive system
would estimate sample difficulty or confidence and allocate extra refinement
only where it is needed. Easy samples could retain one-step or low-step
efficiency, while difficult samples could receive additional corrections.
Candidate stopping signals include the norm of the predicted residual
displacement, disagreement between consecutive predictions, an estimate of
local truncation error, or an auxiliary confidence head. Evaluation should
compare FID, IS, average NFE, wall-clock latency, and the distribution of per-
sample computation. A practical batched implementation may use a small set of
discrete budgets, such as 1, 2, 4, and 8 NFE, to limit branch divergence.

\subsection{Multi-Scale MeanFlow for High-Resolution Generation}

A coarse-to-fine Mean Flow system would first generate global structure at low
resolution and then refine it progressively at higher resolutions. This avoids
performing every operation at full spatial cost. One implementation is a cascade
in which a low-resolution sample is upsampled and supplied as conditioning to a
higher-resolution stage; another is a shared network with resolution-specific
heads.

The central research question is whether Mean Flow's low-NFE advantage survives
across scales. Each stage must remain consistent with the preceding resolution,
and evaluation must report end-to-end network evaluations, latency, memory, and
quality rather than assessing each stage in isolation.
